\section{Lower bounds from shared pilot lists}
\label{sec:pilot-lower}

We use the half-\(\ell_1\) convention for total variation. Two laws
on the space of ordered records with TV less than \(1/3\) cannot have
rejection probabilities at most \(1/3\) and at least \(2/3\), respectively.
This remains true when the second law mixes product laws of strictly
far alternatives, since the power bound then holds under the mixture.

\paragraph{A lower bound near the maximum distance.}
Put \(\delta=1-e>0\) and
\[
 J_\delta(q)=\sum_i q_i\min\{1,\sqrt{q_i/\delta}\}.
\]
We prove that a successful positive budget \(n\) satisfies
\(nJ_\delta(q)>1/64\). Suppose instead \(nJ_\delta(q)\leq 1/64\).
Remove the atoms with \(q_i>\delta/16\). Their total mass \(R\) is
at most \(4J_\delta(q)\), so \(R\leq 1/16\). Write \(q^\circ\) for
the remaining reference normalized by \(1-R\).
The original and normalized ordered product laws differ in TV by at
most \(1-(1-R)^n\leq nR\leq 1/16\). This follows by conditioning the
original record on avoiding all removed labels.

Partition the retained positive atoms into occupied dyadic bins
\((u_j/2,u_j]\) with \(u_j=(\delta/16)2^{-j}\). Let \(b_j\) be the
bin mass under \(q^\circ\), and put \(M=\sum_j b_j\sqrt{u_j}\).
For a retained atom \(q_i\) in bin \(j\), we have
\(\sqrt{u_j}\leq 2\sqrt{q_i}\). Thus
\[
 \sqrt{\delta}\,M
 \leq \frac{2\delta}{1-R}
          \sum_{i:\,q_i\leq \delta/16}q_i\sqrt{q_i/\delta}
 \leq 4\delta J_\delta(q).
\]
We also have \(\sum_j u_j\leq \delta/8\).
For each bin draw independent pilot labels from its conditional
reference, with list length
\[
 k_j=\left\lceil n\sqrt{\delta}\,b_j/\sqrt{u_j}\right\rceil.
\]
A realized alternative allocates mass \(b_j\) to the uniform
distribution on the list slots in bin \(j\), counting repeated
labels with their multiplicities. Each realized law is normalized.
The reference mass of its support is at most
\[
 \sum_j k_j u_j
 \leq n\sqrt{\delta}\,M+\sum_j u_j
 \leq \delta/16+\delta/8<\delta.
\]
Its overlap with \(q\) is therefore strictly less than \(\delta\).

In the ordered predictive law, each coordinate independently
chooses its bin and a uniform slot from the shared list. Couple
this law to \((q^\circ)^{\otimes n}\) by using the corresponding
pilot label on the first use of each slot and fresh independent
conditional draws for repeat uses in the reference record. A
difference can occur only when a pair uses the same bin and slot.
The resulting TV bound is
\[
 \binom n2\sum_j\frac{b_j^2}{k_j}
 \leq \frac{n}{\sqrt{\delta}}M
 \leq 4nJ_\delta(q)\leq 1/16.
\]
Together the two TV costs are at most \(1/8\), which yields the
contradiction. If the strict alternative is nonempty, zero draws
cannot give the two errors, so the bound also provides a positive
sample-scale lower bound in cases where \(1/(64J_\delta(q))<1\).

\paragraph{The regular pilot construction.}
Let \(H\) be the chosen head, \(T\) its tail, \(w=q(T)\), and
\(\Delta=e-w\). For any full law \(p\), define
\[
 L(p)=\sum_{i\in H}(q_i-p_i)_+,\qquad
 O(p)=\sum_{i\in T}\min(q_i,p_i).
\]
The exact identity is \(\operatorname{TV}(p,q)=w+L(p)-O(p)\).
Keep the head labels and replace \(T\) by a single marker, writing
\(\bar q\) for this reduced reference.

Suppose \(w>0\), \(n\geq1\), and \(x>0\) obeys
\[
 \rho=K_T(n,x)
  :=\sum_{i\in T}\min\{q_i/x,\ n q_i^{3/2}/\sqrt{x}\}
  \leq 1/64,\qquad nx\geq 8.
\]
Freeze the tail atoms with \(q_i\geq \tau=1/(n^2x)\).
The two terms in the minimum cross at \(\tau\), giving
\[
 \sum_{i\in T:\,q_i\geq\tau}q_i\leq \rho x,\qquad
 n\sqrt{x}\sum_{i\in T:\,q_i<\tau}q_i^{3/2}\leq \rho x.
\]
Put the other positive atoms in occupied dyadic bins
\((u_j/2,u_j]\) with \(u_j=\tau 2^{-j}\). Denote their
\emph{unconditional} masses under \(q\) by \(b_j\), and set
\(M=\sum_j b_j\sqrt{u_j}\). Then
\[
 n\sqrt{x}\,M\leq 2\rho x,\qquad
 \sum_j u_j\leq 2\tau\leq x/32.
\]
For each bin draw an independent list of
\[
 k_j=\left\lceil n\sqrt{x}\,b_j/\sqrt{u_j}\right\rceil
\]
labels from the conditional reference in that bin.

The channel takes an ordered reduced record. It first draws the lists,
independently of the input record. It keeps every head label. For
each marker it independently chooses a frozen label \(i\) with
weight \(q_i/w\), or a bin \(j\) with weight \(b_j/w\) followed by
a uniform slot in that list. This prescription defines one common
Markov channel on ordered records. Conditional on the lists, an
IID reduced input yields an IID full output. For any reduced input
law and every realized list collection, the full tail overlap is
at most the reference mass of the frozen labels and the occupied
slots. This is bounded by
\[
 \rho x+\sum_j k_j u_j
 \leq 3\rho x+x/32\leq x.
\]
In particular, this bound holds for every possible probability of
the marker in a reduced alternative.

Under \(\bar q^{\otimes n}\), each light bin is selected with its
unconditional probability \(b_j\). The same first-use coupling
works on ordered records, leaving the head and frozen draws
fresh. The predictive TV cost of the fill under this reduced
null is bounded by
\[
 \binom n2\sum_j\frac{b_j^2}{k_j}
  \leq \frac{nM}{\sqrt{x}}\leq 2\rho .
\]
Let a prior on reduced laws have predictive TV at most \(1/8\)
from \(\bar q^{\otimes n}\), and suppose each positive-prior reduced
law has head deficit strictly larger than \(\Delta+x\). The
full filled components all have strict TV greater than \(e\).
By data processing for the common channel and the preceding
null comparison, their predictive TV from \(q^{\otimes n}\) is
at most \(1/8+2\rho<1/3\). No full test with both errors at most
\(1/3\) succeeds.

If \(X(n)\) is the positive root of \(K_T(n,X)=1\), take
\[
 x=4096\max\{X(n),1/n\}.
\]
This meets the two scalar conditions. Each summand shrinks by at
least a factor of \(64\) when the second argument is increased by
a factor of \(4096\). Labels of zero null mass have zero pilot mass. If \(w=0\),
a reduced head prior whose marker remains zero transfers by
keeping its head draws, at zero additional predictive cost. Full
tail overlap is then zero, including when the declared tail is
empty.

\paragraph{The marker fluctuation.}
Write \(A=1-w\). For \(w>0\) and \(nA\geq1\), let
\[
 d=\frac18\sqrt{wA/n}.
\]
Since \(w\leq1\leq nA\), we have \(d\leq A/8\).
Consider the reduced law whose marker has mass \(w+d\) and
whose head probabilities are \(q_i(1-d/A)\). This is a probability
law, with head deficit exactly \(d\). Its one-draw chi-square
divergence with respect to \(\bar q\) is
\[
 d^2/w+d^2/A=d^2/(wA)=1/(64n).
\]
Its ordered product TV from the reduced null is bounded by
\[
 \frac12\sqrt{(1+1/(64n))^n-1}
 \leq \frac{1}{2\sqrt{63}}<\frac1{12}.
\]
The middle inequality follows from the binomial expansion and
\(\binom nk\leq n^k\), followed by the geometric sum.
This singleton reduced prior meets the predictive condition
above. If a full test succeeds, it follows that \(d\leq\Delta+x\),
or
\[
 \sqrt{w(1-w)/n}\leq 8(\Delta+x).
\]
When \(w=0\) the left-hand side vanishes and the same necessity
holds for \(\Delta\geq0\).
